\documentclass{subfiles}
\begin{document}
    Named after \emph{Jaques Hadamard}, was actually proposed by \emph{R. C. Bose} and 
        \emph{S. S. Shrikhande}, and refers to the family of \([2^{k}, k, 2^{k - 1}]\)
        linar code. We can distinguish two ``classes'' of Hadamard codes: 
        the plain one and the augmented one. The former have a message length of 
        \(k = m\); the latter of \(k = m + 1\). In what follows, we refer to the 
        second class of codes as augmented Hadamard codes.

    To construct an Hadamard code, fix \(\gf{2}\) as field and let \(k\) be the 
        lenght of the words. Then, the generator matrix for the code is composed of 
        all binary strings of lenght \(k\), lexicographically ordered. That is,
        the matrix 
        \[G = \begin{pmatrix}
                \uparrow & \uparrow & & \uparrow \\ 
                y_{1} & y_{2} & \cdots & y_{2^{k}} \\ 
                \downarrow & \downarrow & & \downarrow
            \end{pmatrix}\]
        where \(y_{i} \in \set{0, 1}^{k}\).

    With the same assumptions, the generator matrix for an augmented Hadamard code 
        is composed by all the binary strings of lenght \(k\) whose first component,
        is a 1.
\end{document}
